\documentclass[aspectratio=169, 11pt]{beamer}
\usepackage{beamer_preamble}

\begin{document}

% ==== Title ===============================================================
\titleslide{Lesson 6-5 \textperiodcentered\ Quick}{Properties of Logarithms + DOK-3 Power Property Proof \& Richter Generalization}

% ==== Do Now ==============================================================
\begin{frame}{Do Now --- Log Properties Warm-Up}{Algebra 2 \textperiodcentered\ Lesson 6-5 \textperiodcentered\ Quick}
\phasebadges{1--2}{7}{bridge to log properties launch}
\vspace{0.2cm}
\begin{projcard}{3 items: error analysis + expand + condense}{slidegold}
\begin{enumerate}
  \setlength{\itemsep}{4pt}
  \item \textbf{Error Analysis} A student approximates \(\log_2 9\) by entering
        \(\dfrac{\ln 2}{\ln 9}\) on her calculator. Why is this wrong?
  \item \textbf{Expand:} \(\log_6\!\dfrac{49}{5}\)
  \item \textbf{Condense:} \(5\ln s + 6\ln t\)
\end{enumerate}
\end{projcard}
\end{frame}

% ==== Launch ==============================================================
\begin{frame}{Launch --- Example 3: Condense into a Single Log}{Algebra 2 \textperiodcentered\ Lesson 6-5 \textperiodcentered\ Quick}
\phasebadges{1--2}{5}{teacher models}
\vspace{0.15cm}
\begin{projcard}{Rules Card (keep visible all period)}{slidegreen}
\small
\textbf{PRODUCT:} \(\log_b(MN) = \log_b M + \log_b N\)\\[2pt]
\textbf{QUOTIENT:} \(\log_b(M/N) = \log_b M - \log_b N\)\\[2pt]
\textbf{POWER:} \(\log_b(M^n) = n\cdot\log_b M\)\\[2pt]
\textbf{CHANGE OF BASE:} \(\log_b x = (\log x)/(\log b)\)

\smallskip
\textcolor{slideaccent}{\textbf{Condensing order: POWER first \(\rightarrow\) then PRODUCT / QUOTIENT}}
\end{projcard}

\vspace{0.1cm}
\begin{projcard}{Teacher models: \(4\log_4 m + 3\log_4 n - \log_4 p\)}{slidegreen}
\small
Step 1 (Power): \(\log_4(m^4) + \log_4(n^3) - \log_4 p\)\\
Step 2 (Product then Quotient): \(\log_4\!\left(\dfrac{m^4 n^3}{p}\right)\)
\end{projcard}
\end{frame}

% ==== Explore =============================================================
\begin{frame}{Explore --- TPS + DOK-3 Paired Driver}{Algebra 2 \textperiodcentered\ Lesson 6-5 \textperiodcentered\ Quick}
\phasebadges{2--3}{28}{student work}
\small\textcolor{slidegray}{6 items in order. Plan \(\sim\)12 min for \#13 + \#40. Wed-F: skip \#3.}

\vspace{0.1cm}
\begin{projcard}{Practice \#19 \textperiodcentered\ Condense}{slideblue}
\(\log_5 6 + \tfrac{1}{2}\log_5 y\)
\end{projcard}
\vspace{-0.12cm}
\begin{projcard}{Practice \#21 \textperiodcentered\ Condense}{slideblue}
\(\tfrac{1}{3}\ln 27 - 3\ln(2y)\)
\end{projcard}
\vspace{-0.12cm}
\begin{projcard}{Practice \#34 \textperiodcentered\ Solve (change of base)}{slideblue}
\(7^x = 100\)
\end{projcard}
\vspace{-0.12cm}
\begin{projcard}{Practice \#3 \textperiodcentered\ Error analysis (skip Wed-F)}{slideblue}
\(\log_4(c^2 d^5) = 2\cdot\log_4 c \cdot 5\cdot\log_4 d\) --- what's wrong?
\end{projcard}
\vspace{-0.12cm}
\begin{projcard}{Practice \#13 + \#40 Part C \textperiodcentered\ DOK-3 PAIRED (\(\bigstar\))}{slideaccent}
\#13: Prove the Power Property. \quad \#40: Richter generalization.
\end{projcard}
\end{frame}

% ==== Share / Summary =====================================================
\begin{frame}{Share / Summary}{Algebra 2 \textperiodcentered\ Lesson 6-5 \textperiodcentered\ Quick}
\phasebadges{2}{5}{DOK-3 reveal + Topic 6 bridge}
\vspace{0.2cm}
\begin{projcard}{Sentence frame \textperiodcentered\ Power Property Proof (\#13)}{slideblue}
\small
``Let \(x = \log_b M\). Then \(b^x = M\). Raise both sides to \(n\):
\((b^x)^n = M^n\), so \(b^{(nx)} = M^n\).
By the definition of log, \(\log_b(M^n) = nx = n\cdot\log_b M\).''
\end{projcard}

\vspace{0.15cm}
\begin{projcard}{Sentence frame \textperiodcentered\ Richter (\#40 Part C)}{slideblue}
\small
\(R(150A_0) - R(A_0) = \log(150A_0/S) - \log(A_0/S) = \log 150 \approx 2.18\).\\
The larger earthquake is about \textbf{2.18 Richter units} greater.
\end{projcard}

\vspace{0.1cm}
\begin{projcard}{Topic 6 Assessment Bridge}{slideaccent}
\small
\#13 proof + \#40 Part C = the exact DOK-3 pattern on Form B.
\textbf{You have now rehearsed both.}
Form B Q8 and Q15 (geometric series) are \textbf{NOT} on the assessment.
\end{projcard}
\end{frame}

% ==== Exit Ticket =========================================================
\begin{frame}{Exit Ticket --- Two-Part Summary}{Algebra 2 \textperiodcentered\ Lesson 6-5 \textperiodcentered\ Quick}
\phasebadges{1--2}{5}{not graded}
\vspace{0.2cm}
\begin{projcard}{Complete on your exit ticket}{slidegold}
A biologist models a population as \(P = 3\log t + \log 4\).

\vspace{0.2cm}
\begin{enumerate}
  \setlength{\itemsep}{0.3cm}
  \item Write \(P\) as a single logarithm: \(P = \log(\underline{\hspace{2cm}})\)
  \item Solve for \(t\) when \(P = 2\): \quad \(t = \underline{\hspace{2cm}}\)
  \item One biggest thing I learned today: \underline{\hspace{5cm}}
\end{enumerate}
\end{projcard}
\end{frame}

\end{document}
